%!TeX root=../houndtop.tex
\chapter{Extract from the Diary of Dr Watson}
\lettrine[lines=4]{S}{o} far I have been able to quote from the reports which I have forwarded during these early days to Sherlock Holmes. Now, however, I have arrived at a point in my narrative where I am compelled to abandon this method and to trust once more to my recollections, aided by the diary which I kept at the time. A few extracts from the latter will carry me on to those scenes which are indelibly fixed in every detail upon my memory. I proceed, then, from the morning which followed our abortive chase of the convict and our other strange experiences upon the moor.
\vfill

\noindent\hfil\rule{0.5\textwidth}{.4pt}\hfil 

\textsc{October 16th.}—A dull and foggy day with a drizzle of rain. The house is banked in with rolling clouds, which rise now and then to show the dreary curves of the moor, with thin, silver veins upon the sides of the hills, and the distant boulders gleaming where the light strikes upon their wet faces. It is melancholy outside and in. The baronet is in a black reaction after the excitements of the night. I am conscious myself of a weight at my heart and a feeling of impending danger—ever present danger, which is the more terrible because I am unable to define it.

And have I not cause for such a feeling? Consider the long sequence of incidents which have all pointed to some sinister influence which is at work around us. There is the death of the last occupant of the Hall, fulfilling so exactly the conditions of the family legend, and there are the repeated reports from peasants of the appearance of a strange creature upon the moor. Twice I have with my own ears heard the sound which resembled the distant baying of a hound. It is incredible, impossible, that it should really be outside the ordinary laws of nature. A spectral hound which leaves material footmarks and fills the air with its howling is surely not to be thought of. Stapleton may fall in with such a superstition, and Mortimer also; but if I have one quality upon earth it is common-sense, and nothing will persuade me to believe in such a thing. To do so would be to descend to the level of these poor peasants, who are not content with a mere fiend dog but must needs describe him with hell-fire shooting from his mouth and eyes. Holmes would not listen to such fancies, and I am his agent. But facts are facts, and I have twice heard this crying upon the moor. Suppose that there were really some huge hound loose upon it; that would go far to explain everything. But where could such a hound lie concealed, where did it get its food, where did it come from, how was it that no one saw it by day? It must be confessed that the natural explanation offers almost as many difficulties as the other. And always, apart from the hound, there is the fact of the human agency in London, the man in the cab, and the letter which warned Sir Henry against the moor. This at least was real, but it might have been the work of a protecting friend as easily as of an enemy. Where is that friend or enemy now? Has he remained in London, or has he followed us down here? Could he—could he be the stranger whom I saw upon the tor?

It is true that I have had only the one glance at him, and yet there are some things to which I am ready to swear. He is no one whom I have seen down here, and I have now met all the neighbours. The figure was far taller than that of Stapleton, far thinner than that of Frankland. Barrymore it might possibly have been, but we had left him behind us, and I am certain that he could not have followed us. A stranger then is still dogging us, just as a stranger dogged us in London. We have never shaken him off. If I could lay my hands upon that man, then at last we might find ourselves at the end of all our difficulties. To this one purpose I must now devote all my energies.


%\makeatletter
%\@ifclasswith{scrbook}{a5paper}
%{%
%	\begin{wrapfigure}[18]{O}{0.6\textwidth}
%		\centering
%		\includegraphics[width=.7\linewidth]{10_butlerstand}
%		\caption{The butler was standing very pale but very collected before us.}
%	\end{wrapfigure}
%}{%
%}
%\makeatother

My first impulse was to tell Sir Henry all my plans. My second and wisest one is to play my own game and speak as little as possible to anyone. He is silent and distrait. His nerves have been strangely shaken by that sound upon the moor. I will say nothing to add to his anxieties, but I will take my own steps to attain my own end.

We had a small scene this morning after breakfast. Barrymore asked leave to speak with Sir Henry, and they were closeted in his study some little time. Sitting in the billiard-room I more than once heard the sound of voices raised, and I had a pretty good idea what the point was which was under discussion. After a time the baronet opened his door and called for me.

»Barrymore considers that he has a grievance,« he said. »He thinks that it was unfair on our part to hunt his brother-in-law down when he, of his own free will, had told us the secret.«

The butler was standing very pale but very collected before us.

»I may have spoken too warmly, sir,« said he, »and if I have, I am sure that I beg your pardon. At the same time, I was very much surprised when I heard you two gentlemen come back this morning and learned that you had been chasing Selden. The poor fellow has enough to fight against without my putting more upon his track.«

»If you had told us of your own free will it would have been a different thing,« said the baronet, »you only told us, or rather your wife only told us, when it was forced from you and you could not help yourself.«

»I didn't think you would have taken advantage of it, Sir Henry—indeed I didn't.«

»The man is a public danger. There are lonely houses scattered over the moor, and he is a fellow who would stick at nothing. You only want to get a glimpse of his face to see that. Look at Mr Stapleton's house, for example, with no one but himself to defend it. There's no safety for anyone until he is under lock and key.«



»He'll break into no house, sir. I give you my solemn word upon that. But he will never trouble anyone in this country again. I assure you, Sir Henry, that in a very few days the necessary arrangements will have been made and he will be on his way to South America. For God's sake, sir, I beg of you not to let the police know that he is still on the moor. They have given up the chase there, and he can lie quiet until the ship is ready for him. You can't tell on him without getting my wife and me into trouble. I beg you, sir, to say nothing to the police.«



»What do you say, Watson?«

I shrugged my shoulders. »If he were safely out of the country it would relieve the tax-payer of a burden.«

»But how about the chance of his holding someone up before he goes?«

»He would not do anything so mad, sir. We have provided him with all that he can want. To commit a crime would be to show where he was hiding.«

»That is true,« said Sir Henry. »Well, Barrymore\longdash«

	\begin{figure}[tbh]
		\centering
		\includegraphics[width=.7\linewidth]{10_butlerstand}
		\caption{The butler was standing very pale but very collected before us.}
	\end{figure}

»God bless you, sir, and thank you from my heart! It would have killed my poor wife had he been taken again.«

»I guess we are aiding and abetting a felony, Watson? But, after what we have heard I don't feel as if I could give the man up, so there is an end of it. All right, Barrymore, you can go.«

With a few broken words of gratitude the man turned, but he hesitated and then came back.

»You've been so kind to us, sir, that I should like to do the best I can for you in return. I know something, Sir Henry, and perhaps I should have said it before, but it was long after the inquest that I found it out. I've never breathed a word about it yet to mortal man. It's about poor Sir Charles's death.«

The baronet and I were both upon our feet. »Do you know how he died?«

»No, sir, I don't know that.«

»What then?«

»I know why he was at the gate at that hour. It was to meet a woman.«

»To meet a woman! He?«

»Yes, sir.«

»And the woman's name?«

»I can't give you the name, sir, but I can give you the initials. Her initials were L. L.«

»How do you know this, Barrymore?«

»Well, Sir Henry, your uncle had a letter that morning. He had usually a great many letters, for he was a public man and well known for his kind heart, so that everyone who was in trouble was glad to turn to him. But that morning, as it chanced, there was only this one letter, so I took the more notice of it. It was from Coombe Tracey, and it was addressed in a woman's hand.«

»Well?«


»Well, sir, I thought no more of the matter, and never would have done had it not been for my wife. Only a few weeks ago she was cleaning out Sir Charles's study—it had never been touched since his death—and she found the ashes of a burned letter in the back of the grate. The greater part of it was charred to pieces, but one little slip, the end of a page, hung together, and the writing could still be read, though it was gray on a black ground. It seemed to us to be a postscript at the end of the letter, and it said:»Please, please, as you are a gentleman, burn this letter, and be at the gate by ten o'clock.« Beneath it were signed the initials L. L.«

»Have you got that slip?«

»No, sir, it crumbled all to bits after we moved it.«

»Had Sir Charles received any other letters in the same writing?«

»Well, sir, I took no particular notice of his letters. I should not have noticed this one, only it happened to come alone.«

»And you have no idea who L. L. is?«

»No, sir. No more than you have. But I expect if we could lay our hands upon that lady we should know more about Sir Charles's death.«

»I cannot understand, Barrymore, how you came to conceal this important information.«

»Well, sir, it was immediately after that our own trouble came to us. And then again, sir, we were both of us very fond of Sir Charles, as we well might be considering all that he has done for us. To rake this up couldn't help our poor master, and it's well to go carefully when there's a lady in the case. Even the best of us\longdash«

»You thought it might injure his reputation?«

»Well, sir, I thought no good could come of it. But now you have been kind to us, and I feel as if it would be treating you unfairly not to tell you all that I know about the matter.«

»Very good, Barrymore; you can go.« When the butler had left us Sir Henry turned to me. »Well, Watson, what do you think of this new light?«

»It seems to leave the darkness rather blacker than before.«

»So I think. But if we can only trace L. L. it should clear up the whole business. We have gained that much. We know that there is someone who has the facts if we can only find her. What do you think we should do?«

»Let Holmes know all about it at once. It will give him the clue for which he has been seeking. I am much mistaken if it does not bring him down.«

I went at once to my room and drew up my report of the morning's conversation for Holmes. It was evident to me that he had been very busy of late, for the notes which I had from Baker Street were few and short, with no comments upon the information which I had supplied and hardly any reference to my mission. No doubt his blackmailing case is absorbing all his faculties. And yet this new factor must surely arrest his attention and renew his interest. I wish that he were here.

\noindent\hfil\rule{0.5\textwidth}{.4pt}\hfil 

\textsc{October 17th.}—All day to-day the rain poured down, rustling on the ivy and dripping from the eaves. I thought of the convict out upon the bleak, cold, shelterless moor. Poor devil! Whatever his crimes, he has suffered something to atone for them. And then I thought of that other one—the face in the cab, the figure against the moon. Was he also out in that deluged—the unseen watcher, the man of darkness? In the evening I put on my waterproof and I walked far upon the sodden moor, full of dark imaginings, the rain beating upon my face and the wind whistling about my ears. God help those who wander into the great mire now, for even the firm uplands are becoming a morass. I found the black tor upon which I had seen the solitary watcher, and from its craggy summit I looked out myself across the melancholy downs. Rain squalls drifted across their russet face, and the heavy, slate-coloured clouds hung low over the landscape, trailing in gray wreaths down the sides of the fantastic hills. In the distant hollow on the left, half hidden by the mist, the two thin towers of Baskerville Hall rose above the trees. They were the only signs of human life which I could see, save only those prehistoric huts which lay thickly upon the slopes of the hills. Nowhere was there any trace of that lonely man whom I had seen on the same spot two nights before.

\begin{figure}[p]
\centering
\includegraphics[width=1\linewidth]{craggy}
\caption{From its craggy summit I looked out myself across the melancholy downs.}
\end{figure}

As I walked back I was overtaken by Dr Mortimer driving in his dog-cart over a rough moorland track which led from the outlying farmhouse of Foulmire. He has been very attentive to us, and hardly a day has passed that he has not called at the Hall to see how we were getting on. He insisted upon my climbing into his dog-cart, and he gave me a lift homeward. I found him much troubled over the disappearance of his little spaniel. It had wandered on to the moor and had never come back. I gave him such consolation as I might, but I thought of the pony on the Grimpen Mire, and I do not fancy that he will see his little dog again.

»By the way, Mortimer,« said I as we jolted along the rough road, »I suppose there are few people living within driving distance of this whom you do not know?«

»Hardly any, I think.«

»Can you, then, tell me the name of any woman whose initials are L. L.?«

He thought for a few minutes.

»No,« said he. »There are a few gipsies and labouring folk for whom I can't answer, but among the farmers or gentry there is no one whose initials are those. Wait a bit though,« he added after a pause. »There is Laura Lyons—her initials are L. L.—but she lives in Coombe Tracey.«

»Who is she?« I asked.

»She is Frankland's daughter.«

»What! Old Frankland the crank?«

»Exactly. She married an artist named Lyons, who came sketching on the moor. He proved to be a blackguard and deserted her. The fault from what I hear may not have been entirely on one side. Her father refused to have anything to do with her because she had married without his consent, and perhaps for one or two other reasons as well. So, between the old sinner and the young one the girl has had a pretty bad time.«

»How does she live?«

»I fancy old Frankland allows her a pittance, but it cannot be more, for his own affairs are considerably involved. Whatever she may have deserved one could not allow her to go hopelessly to the bad. Her story got about, and several of the people here did something to enable her to earn an honest living. Stapleton did for one, and Sir Charles for another. I gave a trifle myself. It was to set her up in a typewriting business.«

He wanted to know the object of my inquiries, but I managed to satisfy his curiosity without telling him too much, for there is no reason why we should take anyone into our confidence. To-morrow morning I shall find my way to Coombe Tracey, and if I can see this Mrs Laura Lyons, of equivocal reputation, a long step will have been made towards clearing one incident in this chain of mysteries. I am certainly developing the wisdom of the serpent, for when Mortimer pressed his questions to an inconvenient extent I asked him casually to what type Frankland's skull belonged, and so heard nothing but craniology for the rest of our drive. I have not lived for years with Sherlock Holmes for nothing.

I have only one other incident to record upon this tempestuous and melancholy day. This was my conversation with Barrymore just now, which gives me one more strong card which I can play in due time.

Mortimer had stayed to dinner, and he and the baronet played ecarté afterwards. The butler brought me my coffee into the library, and I took the chance to ask him a few questions.

»Well,« said I, »has this precious relation of yours departed, or is he still lurking out yonder?«

»I don't know, sir. I hope to heaven that he has gone, for he has brought nothing but trouble here! I've not heard of him since I left out food for him last, and that was three days ago.«

»Did you see him then?«

»No, sir, but the food was gone when next I went that way.«

»Then he was certainly there?«

»So you would think, sir, unless it was the other man who took it.«

I sat with my coffee-cup halfway to my lips and stared at Barrymore.

»You know that there is another man then?«

»Yes, sir; there is another man upon the moor.«

»Have you seen him?«

»No, sir.«

»How do you know of him then?«

\begin{figure}[tbh]
\centering
\includegraphics[width=.7\linewidth]{09_anotherman}
\caption{»You know that there is another man then?«}
\end{figure}

»Selden told me of him, sir, a week ago or more. He's in hiding, too, but he's not a convict as far as I can make out. I don't like it, Dr Watson—I tell you straight, sir, that I don't like it.« He spoke with a sudden passion of earnestness.

»Now, listen to me, Barrymore! I have no interest in this matter but that of your master. I have come here with no object except to help him. Tell me, frankly, what it is that you don't like.«

Barrymore hesitated for a moment, as if he regretted his outburst, or found it difficult to express his own feelings in words.

»It's all these goings-on, sir,« he cried at last, waving his hand towards the rain-lashed window which faced the moor. »There's foul play somewhere, and there's black villainy brewing, to that I'll swear! Very glad I should be, sir, to see Sir Henry on his way back to London again!«

»But what is it that alarms you?«

»Look at Sir Charles's death! That was bad enough, for all that the coroner said. Look at the noises on the moor at night. There's not a man would cross it after sundown if he was paid for it. Look at this stranger hiding out yonder, and watching and waiting! What's he waiting for? What does it mean? It means no good to anyone of the name of Baskerville, and very glad I shall be to be quit of it all on the day that Sir Henry's new servants are ready to take over the Hall.«

»But about this stranger,« said I. »Can you tell me anything about him? What did Selden say? Did he find out where he hid, or what he was doing?«

»He saw him once or twice, but he is a deep one, and gives nothing away. At first he thought that he was the police, but soon he found that he had some lay of his own. A kind of gentleman he was, as far as he could see, but what he was doing he could not make out.«

»And where did he say that he lived?«

»Among the old houses on the hillside—the stone huts where the old folk used to live.«

»But how about his food?«

»Selden found out that he has got a lad who works for him and brings him all he needs. I dare say he goes to Coombe Tracey for what he wants.«

»Very good, Barrymore. We may talk further of this some other time.« When the butler had gone I walked over to the black window, and I looked through a blurred pane at the driving clouds and at the tossing outline of the wind-swept trees. It is a wild night indoors, and what must it be in a stone hut upon the moor. What passion of hatred can it be which leads a man to lurk in such a place at such a time! And what deep and earnest purpose can he have which calls for such a trial! There, in that hut upon the moor, seems to lie the very centre of that problem which has vexed me so sorely. I swear that another day shall not have passed before I have done all that man can do to reach the heart of the mystery.
